% TODO: Tell that we assume that the user clicks on a node to preview the formula and see the dependencies etc. ==> i.e to answer why we even click on a cell ?
% TODO: Mention the spreadsheet presentation ?
\section{The presentation layer}\label{sec: presentation-layer}
% TODO: move these goals to the visualization layer ?
% TODO: use datatable or worksheet ???
% TODO: explain why only present one formula at a time. 
The representation layers main responsibility are the interactions, for example changing a constant value, modify an edge or highlight cells on hover, and rendering, the dual view and flow graph representation. The system uses a dual view, having on the left side the spreadsheet and on the right side the flow graph. The flow graph aims to represent the hidden dependency graph of the spreadsheet. The flow graph is not static but dynamic and the user can interact and modify the flow graph.

\subsection{Layout strategy}
% TODO: mention Juxtaposition
% TODO: also mention if we overlay it on the datatable the tool will automatically suffer from hidden dependencies and visibilty because the datatable itself suffers from it so by splitting the views we also make it independant of the datatables spatiality.
% TODO: also mention data we explicitly use the mechanism of hyperformula to not make an infinite grid ==> if data is more than the viewport can show ==> minimap to give user some intuition of the size of the spreadsheet
The spatial arrangement of the tool and the spreadsheet has a significant impact on the cognitive load during the spreadsheet comprehension.
\begin{definition}
    \label{def:juxtaposability}
    Juxtaposability is the ability to see any two portions on screen side-by-side at the same time. \cite{green_usability_1996}
\end{definition}
\subsubsection{The split view}
As mentioned in section \ref{sec:built-in-tools}, traditional spreadsheet tools for tracing precedents and dependents use an overlay strategy. While this approach directly integrates visual information into the spreadsheet, it fails when dependencies are complex and span multiple worksheets. Since spreadsheet applications can only display one worksheet at a time, it is difficult to have for example an arrow which is rendered across worksheets. This forces the user to navigate away from the current position in order to understand the cross-worksheet dependencies. 

To address this limitation and achieve juxtaposability between the spreadsheet and the flow graph, the system employs a split view layout. The spreadsheet occupies the left side, while the flow graph appears on the right side. This separation allows the system to render the flow graph independently of the current worksheet-view.

The split view layout introduces the cognitive cost known as split-attention effect. Sweller and Chandler describe the split-attention effect in, The split-attention effect as a factor in the design of instruction \cite{chandler_split-attention_1992}, that Learners are often forced to split their attention between and mentally integrate disparate sources of information. In this case the user has to integrate form the flow-graph to the current worksheet or reversed.

However, the split-attention effect is only a disadvantage when all relevant dependencies are visible within a single worksheet. Spreadsheet with complex formulas frequently contain cross-worksheet dependencies. When dependencies span multiple worksheets, an overlay approach also introduces the split-attention effect. The user must still switch between worksheets to view the remaining dependencies and must mentally integrate the information from one worksheet to another.

\subsubsection{The minimap}
Ragavan et al. mention in, Spreadsheet Comprehension: Guesswork, Giving up and Going Back to the Author \cite{srinivasa_ragavan_spreadsheet_2021}, the problem that it is easy to miss data because of the unbounded grid size of spreadsheets. In order to give the user some sense of the size of the spreadsheet, the system implements a minimap which displays the proportional size of the current viewport to the whole spreadsheet and where the viewport is located.
%TODO: Image ?
% TODO: insert image of the two views

\subsection{Flow graph justification}
The decision to represent the hidden data as a flow graph was driven by the nature of the underlying state of spreadsheet formulas, the dependency graph.
A flow graph allows to uncover this hidden state and enables the user to see the dependencies rather than by reading formulas in the cells and parsing it on their own.
The flow direction is left to right. This should mirror the spatial convention for spreadsheets, where the inputs of a formula are located in columns on the left side of the formula.

\subsection{Interaction design}
\label{sec:interaction-design}
The flow graph serves as a visualisation of a formula's dependency structure and an interface for exploring and manipulating it. The interface operates in two modes, the preview mode and the edit mode, where edit mode extends the available interactions to include modifications to the underlying data. The visual distinction between modes serves as a cue to the user, that changes are now possible. Edit Mode is intentionally easy to enter, prioritizing a low viscosity \cite{green_usability_1996}. 

\subsubsection{Preview Mode}
Preview Mode is the default state of the flow graph and is designed for exploration. Hovering over a reference node highlights the corresponding cell or range in its worksheet. If the cell is already visible in the current viewport, this allows the user to see the dependency at a glance without dirupting their current position in the worksheet. Clicking on a reference node explicitly scrolls to the corresponding worksheet and cell or range, enabling closer inspection when the user requests it. This aims to preserve the user's spatial context in the worksheet and only disrupt it on demand. 

\subsubsection{Edit Mode}
% TODO: Maybe use viscosity to explain that the idea behind all these entry actions is that it should be very near to edit changes ==> instead of tracing the whole formula back to than edit it in the datatable.
% TODO: And also the bidirectional interaction ==> closeness of mapping by existing knowledge to drag or select cells should assist user in editing nodes. But also mention because flow-graph complete new environment ==> new interactions to learn therefore even more important to have these bidirectional interactions which are known
Edit Mode extends the interactions available in Preview Mode to allow direct manipulation of the flow graph and its underlying data. It is designed to minimize the distance between recognizing a value to change and modifying it. Rather than requiring the user to trace the node back to its source in a worksheet and modify it there, changes can be initiated directly from the point of interest in the flow graph. This aims to reduce the viscosity \cite{green_usability_1996}.
Edit Mode is always entered through an explicit user action:
\begin{enumerate}
    \item clicking on the edit button
    \item double-clicking on an inline value or address
    \item double-clicking an edge
    \item dragging from a connection handle
\end{enumerate}
The entry action determines the initial editing context. Double-clicking on a node's address allows the user to select a new address directly from the worksheet. This aims to resemble similar actions in other spreadsheet programs where the user can replace addresses by selecting a cell in a worksheet.%TODO: mention closeness of mapping
Double-clicking an inline value transforms the field into an editable input, enabling the user to modify the constant in place. In both cases, the change is committed by pressing Enter or clicking the confirmation button.
Double-clicking an edge enters Edit mode with that edge pre-selected for further manipulations. Dragging from a connection handle initiates edge creation or modification. The user drags from a source handle to a target handle. While dragging the edge invalid handles are grayed out. If the target handle has no existing connection, the new edge connects directly. Else the target handle already has a connection, a popover asks the user to either replace the existing edge or swap positions with it. Clicking only the edit button enters the edit mode without an initial editing context. Once in edit mode, the user can change context freely.
% picture of the component and arrows what i.e value / address / connection handle is.

